\section{小\personMorse{}函数的\personFloer{}复形}
\label{sec:floer-complex-small-morse-function}

由于 $HF_*(H,J)$ 在改变 $(H,J)$ 时彼此同构，只需选择一个特殊对子，即可确定它作为
$\Lambda_\omega$-模的同构类型。在 (Fl89a)、(Fl89b) 中，\personFloer{}对\personMorse{}函数 $f$
和不依赖时间的殆复结构 $J_0$ 使用了特殊情形
\[
 H=\epsilon f,\qquad J\equiv J_0,
\]
其中 $\epsilon>0$ 充分小。本节分析\personFloer{}复形
\[
 \bigl(CF(\epsilon f),\partial_{(\epsilon f,J_0)}\bigr)
\]
的结构。

下述定理实质上由\personFloer{}在 (Fl89b) 中证明。半正情形参见 (HS95)，一般情形的完整
细节参见 (FOn99)、(LT98)。

\begin{theorem}
\label{thm:small-morse-hamiltonian-floer-complex}
设 $f$ 是 $M$ 上任意小\personMorse{}函数，$J_0$ 是相容殆复结构，并且 $f$ 关于度量
$g_{J_0}$ 是\personMorse{}--\personSmale{}的。则存在 $\epsilon_0>0$，使得对所有
$0<\epsilon\leq\epsilon_0$ 都有链同构
\[
 \bigl(CF_*(\epsilon f),\partial_{(\epsilon f,J_0)}\bigr)
 \cong
 \bigl(CM_*(-\epsilon f),
 \partial_{(-\epsilon f,g_{J_0})}^{\mathrm{Morse}}\bigr)
 \otimes\Lambda_\omega^\downarrow.
\]
\end{theorem}

\begin{proof}
对这个特殊对子 $(\epsilon f,J_0)$，每条满足
\[
 \dot\chi-\operatorname{grad}_g(\epsilon f)=0
\]
的梯度流线 $\chi$，都给出方程
\[
 \frac{\partial u}{\partial\tau}
 +J_0\left(\frac{\partial u}{\partial t}-X_{\epsilon f}(u)\right)=0
\]
的不依赖 $t$ 的解 $u(\tau,t)\equiv\chi(\tau)$。(FOn99) 的主引理 22.4
（开弦版本亦见 (Oh96b) 的附录）证明了：只要 $\epsilon$ 充分小，所有这些流线都横截。
此外，对通有的不依赖 $t$ 的 $J\equiv J_0$，任意依赖 $t$ 的族都至少以二维族出现。
因此，只有上述不依赖 $t$ 的流线对\personFloer{}边界映射有贡献，证明完成。
\end{proof}

有了此定理，利用 $(M,\omega)$ 的典范定向应用\personPoincare{}对偶，得到
\[
 \bigl(CM^*(-\epsilon f),\delta_{(-\epsilon f,g_{J_0})}^{\mathrm{Morse}}\bigr)
 \cong
 \bigl(CM_{2n-*}(-\epsilon f),
 \partial_{(-\epsilon f,g_{J_0})}^{\mathrm{Morse}}\bigr).
\]
这里，$HF_*$ 的分次 $*$ 表示\personFloer{}闭链 $\alpha$ 的次数，它由其生成元的\personConley{}--\personZehnder{}指标给出。

本书采用 (Oh05d) 中 $CF_*(H)$ 的分次约定：
\begin{equation}
 \deg([z,w])=\mu_H([z,w])
\label{eq:19-hamiltonian-floer-grading-convention}
\end{equation}
对所有 $[z,w]\in\operatorname{Crit}\mathcal A_H$ 成立。这个约定类似于 (Oh99) 中
拉格朗日子流形情形所用的约定。

接着比较这一分次与 $-f$ 的\emph{负梯度}流方程（也就是 $f$ 的\emph{正梯度}流）
\[
 \dot\chi-\operatorname{grad}f(\chi)=0
\]
的\personMorse{}复形的\personMorse{}分次。它对应于作用量泛函 $\mathcal A_{\epsilon f}$ 的\emph{负梯度}
流。在本书约定下，这给出\personMorse{}指标 $\mu_{-\epsilon f}^{\mathrm{Morse}}(p)$ 与\personConley{}--\personZehnder{}指标 $\mu_{\epsilon f}([p,\widehat p])$ 的关系。（参见 (SZ92) 的
引理 7.2，但要留意哈密顿向量场约定不同：该文定义的 $X_H$ 是本书约定中的
$-X_H$。）这一关系为
\[
 \mu_{\epsilon f}([p,\widehat p])
 =\mu_{-\epsilon f}^{\mathrm{Morse}}(p)-n,
\]
等价地，
\[
 \mu_{-\epsilon f}^{\mathrm{Morse}}(p)
 =\mu_{\epsilon f}([p,\widehat p])+n.
\]

\begin{corollary}
\label{cor:small-hamiltonian-floer-cohomology}
作为 $\mathbb Q$-向量空间或 $\Lambda_\omega^{(0)}$-模，有自然的典范分次同构
\[
 H^*(M)\otimes\Lambda_\omega
 \cong HF_*(\epsilon f,J_0)=:HF_*(\epsilon f).
\]
\end{corollary}

还回忆，若把量子上同调 $QH^*(M)$ 建模为量子链复形 $CQ_*(-\epsilon f)$ 所对应的\personMorse{}复形
\[
 CM_*(-\epsilon f)\otimes\Lambda^\downarrow
\]
的上同调，则作为 $\Lambda_\omega$-模，
$H^*(M)\otimes\Lambda_\omega$ 与 $QH^*(M)$ 同构。于是，作为分次
$\mathbb Q$-向量空间，有如下保持分次的同构：
\[
 QH^{n-k}(M)\longrightarrow QH_{n+k}(M)
 \cong HQ_{n+k}(-\epsilon f)\longrightarrow HF_k(\epsilon f,J_0)
 \longrightarrow HF_k(H,J).
\]
